\chapter{Two Variables: Tables, Scatter Plots, and Correlation}
\label{ch:two-variables}

Chapters 1--3 looked at one variable at a time: its center, its spread, its
quantiles, its shape. Real questions almost always involve two. Does latency grow
with payload size? Do customers who contact support by phone get answers faster?
Does one treatment work better than another? This chapter builds the tools for
describing how two variables move together, and then spends a long time on the
ways those tools mislead.

There are two cases. If both variables are categories (device type, yes/no), we
count pairs in a table. If both are numbers (size, latency), we draw a scatter
plot and summarize it with the correlation coefficient.

\section{Two categorical variables: contingency tables}
\label{sec:corr-tables}

A support team logs 600 tickets. For each ticket it records the channel the
customer used (chat, email, phone) and how long resolution took (same day, 1--3
days, longer). Counting every combination gives a \textbf{contingency table}: a
grid with one row per category of the first variable, one column per category of
the second, and a count in each cell.

\begin{center}
\begin{tabular}{lrrrr}
\toprule
 & Same day & 1--3 days & Longer & Row total \\
\midrule
Chat  & 120 &  60 & 20 & 200 \\
Email &  90 & 150 & 60 & 300 \\
Phone &  70 &  20 & 10 & 100 \\
\midrule
Column total & 280 & 230 & 90 & 600 \\
\bottomrule
\end{tabular}
\end{center}

The row and column totals are the \textbf{marginal distributions}: the
distribution of one variable alone, ignoring the other. The name comes from where
they sit, in the margins of the table. Overall, $280/600 \approx 46.7\%$ of tickets
were resolved the same day.

Raw counts are hard to compare because the rows have different sizes. Email has
90 same-day tickets and phone has 70, but email had three times as many tickets.
The fix is to divide each cell by its row total. The result is a
\textbf{conditional distribution}: the distribution of resolution time
\emph{given} a particular channel.

\begin{center}
\begin{tabular}{lrrr}
\toprule
 & Same day & 1--3 days & Longer \\
\midrule
Chat  & $120/200 = 60\%$ & $60/200 = 30\%$ & $20/200 = 10\%$ \\
Email & $90/300 = 30\%$ & $150/300 = 50\%$ & $60/300 = 20\%$ \\
Phone & $70/100 = 70\%$ & $20/100 = 20\%$ & $10/100 = 10\%$ \\
\midrule
All   & $280/600 \approx 46.7\%$ & $230/600 \approx 38.3\%$ & $90/600 = 15\%$ \\
\bottomrule
\end{tabular}
\end{center}

Now the story is visible. Phone resolves 70\% of tickets the same day; email only
30\%. A \textbf{segmented bar chart} draws each row of this table as a bar of
height 100\%, cut into pieces whose lengths are the conditional proportions.

\begin{figure}[H]
\centering
\begin{tikzpicture}
\begin{axis}[
  xbar stacked, width=11cm, height=5.2cm, bar width=0.55cm,
  xmin=0, xmax=100, xlabel={percent of tickets in the row},
  symbolic y coords={All,Phone,Email,Chat}, ytick=data,
  y tick label style={font=\small},
  legend style={at={(0.5,1.03)}, anchor=south, legend columns=3, font=\small, draw=none},
  legend cell align=left,
  nodes near coords, nodes near coords style={font=\scriptsize, color=white},
  every node near coord/.append style={xshift=0pt},
  point meta=explicit symbolic,
  enlarge y limits=0.18
]
\addplot[fill=targetgreen, draw=white] coordinates {(60,Chat)[60\%] (30,Email)[30\%] (70,Phone)[70\%] (46.7,All)[46.7\%]};
\addplot[fill=sourceblue, draw=white] coordinates {(30,Chat)[30\%] (50,Email)[50\%] (20,Phone)[20\%] (38.3,All)[38.3\%]};
\addplot[fill=trapred, draw=white] coordinates {(10,Chat)[10\%] (20,Email)[20\%] (10,Phone)[10\%] (15,All)[15\%]};
\legend{Same day, 1--3 days, Longer}
\draw[densely dashed, black, thick] (axis cs:46.7,{[normalized]-0.5}) -- (axis cs:46.7,{[normalized]3.5});
\end{axis}
\end{tikzpicture}
\caption{Segmented bar chart of resolution time given channel. The dashed line
marks the overall same-day share, 46.7\%. If channel and resolution time were
unrelated, every green segment would end at that line. Instead email's stops well
short (30\%) and phone's runs well past it (70\%).}
\label{fig:two-variables-segbar}
\end{figure}

\begin{definition}[Association for categorical variables]
Two categorical variables are \textbf{associated} in a data set if the conditional
distribution of one differs across the categories of the other. They show
\textbf{no association} if every row has the same conditional distribution, which
is then equal to the marginal distribution.
\end{definition}

No association has a concrete numerical meaning. If channel told us nothing about
resolution time, then 46.7\% of email's 300 tickets would be same-day, that is
\[
  300 \times \frac{280}{600} = 140 \text{ tickets}.
\]
In general the count we would expect in row $i$, column $j$ under no association is
\[
  \text{expected}_{ij} = \frac{(\text{row } i \text{ total})\times(\text{column } j \text{ total})}{n}.
\]
Email actually has 90 same-day tickets, 50 fewer than expected. Chapter 10 turns
the gap between observed and expected counts into the chi-square test, which asks
whether gaps this large could come from chance alone.

\begin{trapbox}
Conditioning in the other direction answers a different question. Dividing by
\emph{column} totals gives ``of the same-day tickets, what fraction came by
chat?'': $120/280 \approx 42.9\%$. That is not the same as ``of chat tickets,
what fraction were same-day?'': $120/200 = 60\%$. In symbols, $P(\text{chat} \given
\text{same day}) \ne P(\text{same day} \given \text{chat})$. Mixing these two up is
the root of the base-rate fallacy (Probability From the Ground Up covers it with
Bayes' rule). Always say which variable you conditioned on.
\end{trapbox}

\begin{lstlisting}[language=Python]
import pandas as pd
df = pd.DataFrame({"channel": ["Chat"]*200 + ["Email"]*300 + ["Phone"]*100,
                   "res": ["Same"]*120 + ["1-3"]*60 + ["Long"]*20
                        + ["Same"]*90 + ["1-3"]*150 + ["Long"]*60
                        + ["Same"]*70 + ["1-3"]*20 + ["Long"]*10})
pd.crosstab(df.channel, df.res, margins=True)          # counts
pd.crosstab(df.channel, df.res, normalize="index")     # row proportions
\end{lstlisting}

\section{Two numerical variables: scatter plots}
\label{sec:corr-scatter}

When both variables are numbers we plot each observation as a point $(x_i, y_i)$.
This is a \textbf{scatter plot}. By convention the variable we think of as the
input, or the one we could control, goes on the horizontal axis and is called the
\textbf{explanatory variable}; the one we want to understand goes on the
vertical axis and is called the \textbf{response}.

Read a scatter plot in four passes:
\begin{itemize}
  \item \textbf{Direction.} Do points drift up to the right (positive association) or
    down to the right (negative association)?
  \item \textbf{Form.} Do they follow a straight line, a curve, or separate clusters?
  \item \textbf{Strength.} How tightly do they hug that form?
  \item \textbf{Outliers.} Are there points far from the overall pattern?
\end{itemize}

\begin{figure}[H]
\centering
\begin{tikzpicture}
\pgfplotsset{gal/.style={scale only axis, width=4.4cm, height=2.9cm,
  xmin=-0.5, xmax=10.5, ticks=none, title style={font=\small},
  every axis plot/.append style={only marks, mark=*, mark size=1.3pt}}}
\begin{axis}[gal, at={(0,0)}, title={positive, linear, strong}, ymin=0, ymax=11]
\addplot[sourceblue] coordinates {(0.00,2.00) (0.42,2.54) (0.83,2.47) (1.25,2.38) (1.67,3.02) (2.08,2.97) (2.50,4.04) (2.92,5.27) (3.33,4.32) (3.75,4.57) (4.17,5.68) (4.58,5.92) (5.00,6.07) (5.42,5.68) (5.83,6.65) (6.25,7.49) (6.67,6.39) (7.08,7.35) (7.50,6.67) (7.92,7.43) (8.33,7.38) (8.75,8.84) (9.17,8.45) (9.58,9.86) (10.00,10.11)};
\end{axis}
\begin{axis}[gal, at={(6.2cm,0)}, title={negative, linear, weak}, ymin=-1, ymax=13]
\addplot[sourceblue] coordinates {(0.00,8.63) (0.42,3.76) (0.83,7.51) (1.25,8.28) (1.67,8.39) (2.08,4.90) (2.50,6.79) (2.92,5.58) (3.33,5.72) (3.75,9.25) (4.17,5.30) (4.58,6.64) (5.00,8.27) (5.42,5.12) (5.83,5.86) (6.25,6.10) (6.67,5.79) (7.08,3.01) (7.50,5.40) (7.92,7.76) (8.33,1.74) (8.75,6.34) (9.17,4.66) (9.58,2.93) (10.00,8.00)};
\end{axis}
\begin{axis}[gal, at={(0,-4.1cm)}, title={curved, strong, no direction}, ymin=-0.5, ymax=11]
\addplot[sourceblue] coordinates {(0.00,10.13) (0.42,7.75) (0.83,7.11) (1.25,6.21) (1.67,4.79) (2.08,4.32) (2.50,3.15) (2.92,2.85) (3.33,2.69) (3.75,1.21) (4.17,1.34) (4.58,0.83) (5.00,1.06) (5.42,0.47) (5.83,0.95) (6.25,1.45) (6.67,2.42) (7.08,3.09) (7.50,2.53) (7.92,3.58) (8.33,5.21) (8.75,4.93) (9.17,6.84) (9.58,8.30) (10.00,10.38)};
\end{axis}
\begin{axis}[gal, at={(6.2cm,-4.1cm)}, title={positive, linear, one outlier}, ymin=0, ymax=9]
\addplot[sourceblue] coordinates {(0.00,2.41) (0.42,2.05) (0.83,2.28) (1.25,2.60) (1.67,3.91) (2.08,2.99) (2.50,3.32) (2.92,3.96) (3.33,3.93) (3.75,4.13) (4.17,3.83) (4.58,4.74) (5.00,4.73) (5.42,5.95) (5.83,5.89) (6.25,5.74) (6.67,6.40) (7.08,6.05) (7.50,7.13) (7.92,6.75) (8.33,0.50) (8.75,6.48) (9.17,7.71) (9.58,6.74) (10.00,6.78)};
\draw[trapred, thick] (axis cs:8.33,0.5) circle[radius=0.25cm];
\node[font=\scriptsize, trapred, anchor=east] at (axis cs:7.6,0.6) {outlier};
\end{axis}
\end{tikzpicture}
\caption{Four scatter plots, each described by direction, form, strength, and
outliers. The curved panel has a strong relationship but no overall up or down
direction; a single number measuring ``upness'' will miss it. In the last panel
one point (circled) sits far below the line the other 24 follow.}
\label{fig:two-variables-gallery}
\end{figure}

\begin{recoverybox}
Before reading on: in the bottom-right panel, would removing the circled point make
the relationship look stronger or weaker? (Stronger. The point sits far below the
line the other 24 points follow.)
\end{recoverybox}

\section{Covariance: the quadrant picture}
\label{sec:corr-covariance}

We want one number that says ``these points go up together''. Start with six API
requests, recording payload size $x$ in KB and response latency $y$ in ms.

\begin{center}
\begin{tabular}{crrrrr}
\toprule
$i$ & $x_i$ (KB) & $y_i$ (ms) & $x_i - \xbar$ & $y_i - \ybar$ & $(x_i-\xbar)(y_i-\ybar)$ \\
\midrule
1 & 1 & 12 & $-3$ & $-5$ & $15$ \\
2 & 2 & 19 & $-2$ & $2$ & $-4$ \\
3 & 3 & 13 & $-1$ & $-4$ & $4$ \\
4 & 5 & 20 & $1$ & $3$ & $3$ \\
5 & 6 & 22 & $2$ & $5$ & $10$ \\
6 & 7 & 16 & $3$ & $-1$ & $-3$ \\
\midrule
sum & 24 & 102 & 0 & 0 & 25 \\
\bottomrule
\end{tabular}
\end{center}

The means are $\xbar = 24/6 = 4$ KB and $\ybar = 102/6 = 17$ ms. Draw a vertical
line at $\xbar$ and a horizontal line at $\ybar$. They cut the plane into four
quadrants centered on the point of means $(\xbar, \ybar)$.

\begin{figure}[H]
\centering
\begin{tikzpicture}
\begin{axis}[width=10cm, height=7cm, xmin=0, xmax=8, ymin=10, ymax=24,
  xlabel={payload $x$ (KB)}, ylabel={latency $y$ (ms)},
  xtick={1,2,3,4,5,6,7}, ytick={12,14,16,17,18,20,22}]
\fill[targetgreen, opacity=0.18, draw=targetgreen!70!black, draw opacity=1] (axis cs:4,17) rectangle (axis cs:1,12);
\fill[targetgreen, opacity=0.18, draw=targetgreen!70!black, draw opacity=1] (axis cs:4,17) rectangle (axis cs:3,13);
\fill[targetgreen, opacity=0.18, draw=targetgreen!70!black, draw opacity=1] (axis cs:4,17) rectangle (axis cs:5,20);
\fill[targetgreen, opacity=0.18, draw=targetgreen!70!black, draw opacity=1] (axis cs:4,17) rectangle (axis cs:6,22);
\fill[trapred, opacity=0.22, draw=trapred, draw opacity=1] (axis cs:4,17) rectangle (axis cs:2,19);
\fill[trapred, opacity=0.22, draw=trapred, draw opacity=1] (axis cs:4,17) rectangle (axis cs:7,16);
\draw[dashed, gray] (axis cs:4,10) -- (axis cs:4,24);
\draw[dashed, gray] (axis cs:0,17) -- (axis cs:8,17);
\addplot[only marks, mark=*, mark size=2.2pt, targetgreen!60!black] coordinates {(1,12) (3,13) (5,20) (6,22)};
\addplot[only marks, mark=*, mark size=2.2pt, trapred] coordinates {(2,19) (7,16)};
\node[font=\scriptsize, anchor=north] at (axis cs:1,11.85) {area $+15$};
\node[font=\scriptsize, anchor=south east] at (axis cs:2,19) {area $-4$};
\node[font=\scriptsize, anchor=east] at (axis cs:2.9,13) {area $+4$};
\node[font=\scriptsize, anchor=south] at (axis cs:5,20.1) {area $+3$};
\node[font=\scriptsize, anchor=west] at (axis cs:6.05,22) {area $+10$};
\node[font=\scriptsize, anchor=north east] at (axis cs:7,16) {area $-3$};
\node[font=\scriptsize, align=right, anchor=south east] at (axis cs:7.9,17.3) {above both:\\ product $+$};
\node[font=\scriptsize, align=right, anchor=south east] at (axis cs:3.9,10.2) {below both:\\ product $+$};
\node[font=\scriptsize, align=left, anchor=north west] at (axis cs:0.1,23.8) {$x$ below, $y$ above:\\ product $-$};
\node[font=\scriptsize, align=right, anchor=south east] at (axis cs:7.9,10.2) {$x$ above, $y$ below:\\ product $-$};
\node[font=\scriptsize, anchor=south west] at (axis cs:4,17) {$(\xbar,\ybar)$};
\end{axis}
\end{tikzpicture}
\caption{Each point spans a rectangle with the point of means $(\xbar,\ybar) =
(4, 17)$; the label next to each point is the rectangle's signed area
$(x_i-\xbar)(y_i-\ybar)$. Green rectangles (upper right, lower left) count as
positive, red ones as negative. Green wins, $(15+4+3+10) - (4+3) = 25$, so the
association is positive.}
\label{fig:two-variables-rectangles}
\end{figure}

A point in the upper right is above average in both variables, so both deviations
are positive and the product is positive. A point in the lower left is below
average in both, and the product of two negatives is positive again. Points in the
other two quadrants are above average in one variable and below in the other, and
contribute negative products. Adding all the products asks: on balance, do points
fall in the agreeing quadrants or the disagreeing ones?

\begin{definition}[Sample covariance]
For paired data $(x_1,y_1),\dots,(x_n,y_n)$ the \textbf{sample covariance} is
\[
  s_{xy} = \frac{1}{n-1}\sum_{i=1}^{n}(x_i-\xbar)(y_i-\ybar).
\]
It is the average signed rectangle area, with the same $n-1$ divisor as the sample
variance from Chapter 2.
\end{definition}

For the requests, $s_{xy} = 25/5 = 5$ KB$\cdot$ms. Notice that $s_{xx}$, the
covariance of $x$ with itself, is exactly the sample variance $s_x^2$: every
rectangle becomes a square.

With more points the picture is the same. Figure~\ref{fig:two-variables-quadrant}
colours 30 points by the sign of their product.

\begin{figure}[H]
\centering
\begin{tikzpicture}
\begin{axis}[width=10cm, height=6.5cm, xmin=25, xmax=75, ymin=10, ymax=35,
  xlabel={$x$}, ylabel={$y$},
  legend style={at={(0.02,0.98)}, anchor=north west, font=\small}, legend cell align=left]
\addplot[only marks, mark=*, mark size=1.5pt, targetgreen!70!black] coordinates {(26.70,11.20) (63.50,20.80) (52.70,23.20) (49.60,14.80) (39.40,18.30) (61.70,20.20) (37.00,16.30) (48.30,19.00) (62.80,24.10) (61.60,24.40) (39.40,18.60) (69.90,33.20) (61.30,21.60) (59.20,22.40) (43.70,18.60) (47.60,16.90) (65.20,20.10) (42.80,13.80) (54.40,21.90) (45.30,15.30)};
\addplot[only marks, mark=square*, mark size=1.5pt, trapred] coordinates {(44.80,22.90) (36.30,20.10) (46.80,20.10) (51.60,19.00) (48.40,24.00) (47.40,20.50) (46.70,20.90) (44.90,22.00) (41.10,20.50) (60.10,15.40)};
\draw[dashed, gray] (axis cs:50.01,10) -- (axis cs:50.01,35);
\draw[dashed, gray] (axis cs:25,20) -- (axis cs:75,20);
\node[font=\scriptsize, gray!50!black, anchor=south west] at (axis cs:50.3,10.3) {$\xbar = 50$};
\node[font=\scriptsize, gray!50!black, anchor=south east] at (axis cs:74.6,20.1) {$\ybar = 20$};
\legend{product $>0$ (20 points), product $<0$ (10 points)}
\end{axis}
\end{tikzpicture}
\caption{Thirty points split by the dashed lines at $\xbar$ and $\ybar$. Twenty
(green circles) fall in the agreeing quadrants and ten (red squares) in the
disagreeing ones. Most red squares hug the line $y = \ybar$, so their negative
rectangles are thin: the sum of products, and the covariance, is clearly
positive.}
\label{fig:two-variables-quadrant}
\end{figure}

Covariance has a flaw: its size depends on units. Measure payload in bytes instead
of KB and every $x_i - \xbar$ grows by a factor of 1000, so $s_{xy}$ jumps from 5
to 5000. The relationship did not change; only the ruler did. A covariance of 5000
tells you nothing until you know the scales of both variables.

\section{Pearson correlation}
\label{sec:corr-pearson}

The cure for the units problem is the same trick as in Chapter 2: measure each
deviation in standard deviations. Recall that the \textbf{z-score} of $x_i$ is
$z_{x,i} = (x_i - \xbar)/s_x$, the number of standard deviations $x_i$ sits above
its mean. A z-score has no units.

\begin{definition}[Pearson correlation]
The \textbf{(Pearson) correlation coefficient} of paired data is the average product
of z-scores,
\[
  r = \frac{1}{n-1}\sum_{i=1}^{n} z_{x,i}\,z_{y,i}
    = \frac{1}{n-1}\sum_{i=1}^{n}\frac{x_i-\xbar}{s_x}\cdot\frac{y_i-\ybar}{s_y}
    = \frac{s_{xy}}{s_x s_y}.
\]
\end{definition}

So $r$ is covariance computed after rescaling each axis so that one standard
deviation has length 1. The quadrant picture is unchanged: same quadrants, same
signs, only the rectangles are resized.

\begin{example}[Correlation by hand]
For the six requests, $\sum (x_i-\xbar)^2 = 9+4+1+1+4+9 = 28$ and
$\sum (y_i-\ybar)^2 = 25+4+16+9+25+1 = 80$. So
\[
  s_x = \sqrt{28/5} \approx 2.366 \text{ KB}, \qquad s_y = \sqrt{80/5} = 4 \text{ ms}.
\]
Then
\[
  r = \frac{s_{xy}}{s_x s_y} = \frac{5}{2.366 \times 4} \approx 0.528 .
\]
The z-score route gives the same thing. The z-scores are
\begin{align*}
  z_x &: -1.268,\ -0.845,\ -0.423,\ 0.423,\ 0.845,\ 1.268,\\
  z_y &: -1.25,\ 0.5,\ -1,\ 0.75,\ 1.25,\ -0.25.
\end{align*}
The six products are
\[
  1.585,\ -0.423,\ 0.423,\ 0.317,\ 1.056,\ -0.317,
\]
which sum to $2.641$. Dividing by $n-1=5$ gives $r \approx 0.528$. In bytes and
seconds instead of KB and ms, $r$ is still $0.528$.
\end{example}

\begin{figure}[H]
\centering
\begin{tikzpicture}
\begin{axis}[width=7.6cm, height=7.6cm, xmin=-1.7, xmax=1.7, ymin=-1.7, ymax=1.7,
  axis equal image,
  xlabel={payload z-score $z_x$ (SDs from $\xbar$)},
  ylabel={latency z-score $z_y$ (SDs from $\ybar$)},
  xtick={-1,0,1}, ytick={-1,0,1}, grid=major, grid style={gray!20},
  label style={font=\small}, tick label style={font=\small}]
\fill[targetgreen, opacity=0.18, draw=targetgreen!70!black, draw opacity=1] (axis cs:0,0) rectangle (axis cs:-1.268,-1.25);
\fill[targetgreen, opacity=0.18, draw=targetgreen!70!black, draw opacity=1] (axis cs:0,0) rectangle (axis cs:-0.423,-1);
\fill[targetgreen, opacity=0.18, draw=targetgreen!70!black, draw opacity=1] (axis cs:0,0) rectangle (axis cs:0.423,0.75);
\fill[targetgreen, opacity=0.18, draw=targetgreen!70!black, draw opacity=1] (axis cs:0,0) rectangle (axis cs:0.845,1.25);
\fill[trapred, opacity=0.22, draw=trapred, draw opacity=1] (axis cs:0,0) rectangle (axis cs:-0.845,0.5);
\fill[trapred, opacity=0.22, draw=trapred, draw opacity=1] (axis cs:0,0) rectangle (axis cs:1.268,-0.25);
\addplot[only marks, mark=*, mark size=2.2pt, targetgreen!60!black] coordinates {(-1.268,-1.25) (-0.423,-1) (0.423,0.75) (0.845,1.25)};
\addplot[only marks, mark=*, mark size=2.2pt, trapred] coordinates {(-0.845,0.5) (1.268,-0.25)};
\node[font=\scriptsize, anchor=north] at (axis cs:-1.268,-1.30) {$+1.585$};
\node[font=\scriptsize, anchor=north] at (axis cs:-0.423,-1.05) {$+0.423$};
\draw[gray] (axis cs:0.47,0.75) -- (axis cs:1.0,0.75);
\node[font=\scriptsize, anchor=west] at (axis cs:1.0,0.75) {$+0.317$};
\node[font=\scriptsize, anchor=south] at (axis cs:0.845,1.30) {$+1.056$};
\node[font=\scriptsize, anchor=south] at (axis cs:-0.845,0.55) {$-0.423$};
\node[font=\scriptsize, anchor=north] at (axis cs:1.268,-0.30) {$-0.317$};
\end{axis}
\end{tikzpicture}
\caption{Figure~\ref{fig:two-variables-rectangles} redrawn with each axis
measured in standard deviations, so the point of means sits at $(0,0)$ and each
grid square is one SD by one SD. Every point stays in its quadrant with its
colour; only the rectangles are resized, and the number by each point is now its
product of z-scores. The areas add to
$3.381 - 0.740 = 2.641$, and $r$ is that total divided by $n-1=5$: $0.528$.}
\label{fig:two-variables-zscores}
\end{figure}

\begin{lstlisting}[language=Python]
import numpy as np
x = np.array([1, 2, 3, 5, 6, 7]); y = np.array([12, 19, 13, 20, 22, 16])
np.cov(x, y)[0, 1]        # 5.0   (uses n-1)
np.corrcoef(x, y)[0, 1]   # 0.528
\end{lstlisting}

\subsection{Properties of $r$}

\begin{theorem}[Properties of the correlation coefficient]
\label{thm:corr-props}
For any paired data with $s_x > 0$ and $s_y > 0$:
\begin{enumerate}
  \item $-1 \le r \le 1$.
  \item $r = 1$ exactly when all points lie on a line with positive slope, and
    $r=-1$ exactly when they lie on a line with negative slope.
  \item $r$ has no units. Replacing $x$ by $a + bx$ with $b > 0$ (a change of units)
    leaves $r$ unchanged; with $b<0$ it flips the sign of $r$.
  \item $r$ is symmetric: the correlation of $x$ with $y$ equals that of $y$ with $x$.
\end{enumerate}
\end{theorem}

\begin{whybox}
Why can $r$ never pass $\pm 1$? Z-scores satisfy $\sum z_{x,i}^2 = n-1$, because
$\sum (x_i-\xbar)^2 / s_x^2 = (n-1)s_x^2/s_x^2$. The same holds for $y$. A sum of
squares cannot be negative, so
\[
  0 \le \sum_{i=1}^n (z_{x,i} \pm z_{y,i})^2
    = \sum z_{x,i}^2 \pm 2\sum z_{x,i}z_{y,i} + \sum z_{y,i}^2
    = (n-1) \pm 2(n-1)r + (n-1).
\]
Dividing by $2(n-1)$ gives $0 \le 1 \pm r$, so $-1 \le r \le 1$. Equality $r = 1$
needs $\sum(z_{x,i} - z_{y,i})^2 = 0$, so $z_{y,i} = z_{x,i}$ for every point:
$y$ is an exact increasing linear function of $x$. That is property 2. For
property 3, the change $x \mapsto a + bx$ turns the mean into $a + b\xbar$, so each
deviation $x_i - \xbar$ is multiplied by $b$ and $s_x$ by $\abs{b}$; the ratio
$(x_i-\xbar)/s_x$ is multiplied by $b/\abs{b} = \pm 1$.
\end{whybox}

Figure~\ref{fig:two-variables-rgallery} calibrates the eye. Each panel has 40
points with both variables standardized.

\begin{figure}[H]
\centering
\begin{tikzpicture}
\pgfplotsset{rg/.style={scale only axis, width=3.2cm, height=2.6cm,
  xmin=-3.2, xmax=3.2, ymin=-3.2, ymax=3.2, ticks=none, title style={font=\small},
  every axis plot/.append style={only marks, mark=*, mark size=1pt, distblue}}}
\begin{axis}[rg, at={(0,0)}, title={$r=-0.9$}]
\addplot coordinates {(-0.41,0.20) (1.86,-1.23) (-1.38,1.06) (-0.38,0.33) (-1.35,1.14) (0.46,0.31) (0.54,-0.24) (-1.24,1.97) (-1.65,1.77) (-0.10,-0.37) (-0.42,-0.02) (0.60,-0.57) (1.68,-2.22) (0.22,-0.63) (0.99,-1.80) (-0.11,0.14) (-1.23,1.38) (0.40,0.09) (1.17,-1.42) (-0.55,0.31) (0.96,-0.42) (-0.47,0.05) (-0.55,0.13) (-0.01,-0.10) (2.24,-2.07) (-0.06,0.21) (-0.71,0.35) (0.24,0.12) (-1.06,1.43) (-0.12,0.13) (1.06,-0.59) (1.26,-0.98) (0.21,-0.07) (-0.75,1.11) (-0.53,-0.19) (2.24,-1.90) (-0.93,1.54) (-0.10,0.52) (-0.59,0.31) (-1.43,0.22)};
\end{axis}
\begin{axis}[rg, at={(4.2cm,0)}, title={$r=-0.5$}]
\addplot coordinates {(-1.33,0.52) (0.81,1.50) (-1.39,1.47) (-0.88,0.22) (-1.93,-0.34) (-1.09,-0.55) (-0.15,-0.63) (1.36,-0.71) (0.05,0.82) (0.47,-0.40) (-0.26,-0.08) (1.15,-0.06) (0.29,0.40) (-0.81,0.74) (-1.85,2.00) (-0.04,0.06) (0.57,-0.43) (-0.46,1.13) (-0.01,-0.89) (0.38,-1.76) (1.85,-1.24) (-0.09,1.84) (2.44,-1.31) (0.60,0.48) (1.88,-1.22) (0.01,0.91) (0.18,0.43) (0.59,-1.39) (-1.18,-1.72) (-0.46,1.26) (-1.49,1.71) (1.05,-0.72) (-0.16,0.04) (-1.19,0.27) (-0.21,0.51) (0.80,-1.17) (0.85,-1.22) (-0.35,0.38) (-0.18,-0.13) (0.18,-0.68)};
\end{axis}
\begin{axis}[rg, at={(8.4cm,0)}, title={$r=0$}]
\addplot coordinates {(-0.63,-0.30) (0.77,-0.12) (1.58,0.26) (-2.63,-0.65) (1.37,0.03) (-0.44,-0.99) (1.08,1.05) (0.52,0.42) (1.18,0.87) (-1.12,0.69) (0.08,-0.16) (0.60,0.01) (-2.60,1.43) (0.83,-0.13) (-0.38,0.16) (-1.17,0.20) (-0.26,-0.43) (-0.86,-0.81) (0.89,1.38) (0.37,-0.08) (-0.82,0.65) (1.33,-1.66) (0.95,-2.25) (-0.36,0.33) (-1.24,-1.70) (0.10,0.86) (-0.47,1.32) (-0.74,1.73) (0.04,0.19) (0.47,0.31) (1.53,2.33) (0.41,-0.20) (-1.08,-0.90) (-0.05,-0.16) (-0.10,-0.05) (0.37,-1.44) (-0.95,1.06) (0.66,-0.57) (1.06,-0.70) (-0.28,-1.97)};
\end{axis}
\begin{axis}[rg, at={(0,-3.6cm)}, title={$r=0.3$}]
\addplot coordinates {(0.39,0.21) (0.60,1.17) (-0.31,0.14) (0.72,0.17) (-0.54,-0.06) (0.36,1.35) (1.41,0.56) (0.80,-0.37) (0.19,0.33) (1.25,1.36) (-0.79,0.50) (-0.04,-0.28) (-0.64,-1.29) (2.22,1.19) (0.55,-1.07) (-1.38,-0.62) (-0.50,0.90) (-1.30,0.25) (1.06,-0.56) (0.49,-0.84) (-0.03,0.55) (1.27,-0.20) (-0.20,1.32) (-1.46,-0.48) (-1.17,-1.21) (0.28,1.09) (2.06,2.64) (0.85,0.27) (-0.64,0.03) (-2.49,1.13) (-0.99,-1.58) (-0.80,0.86) (0.80,-0.02) (0.04,0.05) (-1.31,-0.22) (-0.35,-1.92) (0.65,-1.55) (-0.16,-1.25) (-1.15,-1.23) (0.28,-1.32)};
\end{axis}
\begin{axis}[rg, at={(4.2cm,-3.6cm)}, title={$r=0.7$}]
\addplot coordinates {(-0.09,1.25) (-0.56,-0.22) (0.71,0.22) (0.53,0.05) (0.74,0.71) (-0.79,0.67) (-0.01,-0.03) (-0.78,-0.69) (-0.01,0.29) (1.11,-0.05) (-0.25,0.88) (0.25,0.64) (1.57,1.46) (0.52,0.48) (1.73,2.08) (-0.60,0.18) (-1.00,0.30) (0.71,0.52) (-0.30,-1.33) (0.79,0.61) (-0.22,0.55) (-0.55,-0.63) (-0.58,-1.20) (2.32,1.56) (-0.26,0.64) (-2.88,-1.77) (0.08,-0.50) (0.49,-0.99) (0.06,0.83) (-2.04,-2.27) (-0.13,0.01) (-0.62,-0.05) (0.20,-0.77) (-0.14,-0.96) (-2.01,-3.03) (0.88,0.54) (1.75,0.13) (-0.42,0.00) (-0.23,0.25) (0.04,-0.38)};
\end{axis}
\begin{axis}[rg, at={(8.4cm,-3.6cm)}, title={$r=0.95$}]
\addplot coordinates {(-2.26,-2.40) (1.57,1.67) (1.25,0.95) (-0.22,-0.50) (0.34,0.32) (0.40,0.93) (-0.81,-0.51) (-1.34,-0.88) (-0.43,-0.43) (0.11,0.63) (-0.04,-0.21) (-1.13,-0.96) (-0.82,-0.92) (-0.04,0.17) (-0.69,-0.42) (1.66,1.04) (1.27,1.85) (-0.57,-0.33) (-1.31,-1.10) (-0.32,-0.20) (-0.88,-1.08) (-0.44,-0.64) (-0.61,-0.63) (-1.08,-1.05) (0.42,-0.37) (0.15,-0.12) (-1.10,-1.46) (-0.85,-1.04) (0.27,-0.17) (0.93,0.57) (0.59,0.87) (-1.20,-0.78) (2.03,1.71) (0.54,0.45) (0.09,0.02) (-0.45,-0.17) (0.87,0.82) (1.48,1.73) (0.73,0.46) (1.87,2.16)};
\end{axis}
\end{tikzpicture}
\caption{Scatter plots with correlations from $-0.9$ to $0.95$, every panel drawn
on the same standardized scale so only the correlation changes. Moderate values
look weaker than you might expect: at $r=0.3$ the cloud is almost round, and only
near $\pm 0.9$ does it look like a line.}
\label{fig:two-variables-rgallery}
\end{figure}

\begin{trapbox}
$r$ measures \emph{linear} association only, and only its strength, not its slope.
A cloud with $r=0.95$ can have a slope of $0.001$ or of $1000$; the slope depends on
units and $r$ does not. ``Strong correlation'' does not mean ``big effect''.
\end{trapbox}

\section{Always plot: Anscombe's quartet}
\label{sec:corr-anscombe}

In 1973 the statistician Francis Anscombe published four small data sets, each with
11 points, built to make one point. Here they are in full.

\begin{center}
\small
\begin{tabular}{rr|rr|rr|rr}
\toprule
\multicolumn{2}{c|}{I} & \multicolumn{2}{c|}{II} & \multicolumn{2}{c|}{III} & \multicolumn{2}{c}{IV} \\
$x$ & $y$ & $x$ & $y$ & $x$ & $y$ & $x$ & $y$ \\
\midrule
10 & 8.04 & 10 & 9.14 & 10 & 7.46 & 8 & 6.58 \\
8 & 6.95 & 8 & 8.14 & 8 & 6.77 & 8 & 5.76 \\
13 & 7.58 & 13 & 8.74 & 13 & 12.74 & 8 & 7.71 \\
9 & 8.81 & 9 & 8.77 & 9 & 7.11 & 8 & 8.84 \\
11 & 8.33 & 11 & 9.26 & 11 & 7.81 & 8 & 8.47 \\
14 & 9.96 & 14 & 8.10 & 14 & 8.84 & 8 & 7.04 \\
6 & 7.24 & 6 & 6.13 & 6 & 6.08 & 8 & 5.25 \\
4 & 4.26 & 4 & 3.10 & 4 & 5.39 & 19 & 12.50 \\
12 & 10.84 & 12 & 9.13 & 12 & 8.15 & 8 & 5.56 \\
7 & 4.82 & 7 & 7.26 & 7 & 6.42 & 8 & 7.91 \\
5 & 5.68 & 5 & 4.74 & 5 & 5.73 & 8 & 6.89 \\
\bottomrule
\end{tabular}
\end{center}

All four have $\xbar = 9$, $s_x^2 = 11$, $\ybar \approx 7.50$, $s_y^2 \approx
4.12$--$4.13$, $r \approx 0.816$, and the same best-fit line $\hat y \approx 3.00 + 0.500x$
(Chapter 12 explains how that line is found). A table of summary statistics
cannot tell them apart. A picture can.

\begin{figure}[H]
\centering
\begin{tikzpicture}
\pgfplotsset{an/.style={scale only axis, width=4.6cm, height=3.3cm,
  xmin=2, xmax=20, ymin=2, ymax=14, xtick={4,8,12,16,20}, ytick={4,8,12},
  tick label style={font=\scriptsize}, title style={font=\small}}}
\begin{axis}[an, at={(0,0)}, title={I: linear with scatter}]
\addplot[only marks, mark=*, mark size=1.3pt, sourceblue] coordinates {(10,8.04) (8,6.95) (13,7.58) (9,8.81) (11,8.33) (14,9.96) (6,7.24) (4,4.26) (12,10.84) (7,4.82) (5,5.68)};
\addplot[densityred, thick, domain=2:20] {3+0.5*x};
\end{axis}
\begin{axis}[an, at={(6.3cm,0)}, title={II: a smooth curve}]
\addplot[only marks, mark=*, mark size=1.3pt, sourceblue] coordinates {(10,9.14) (8,8.14) (13,8.74) (9,8.77) (11,9.26) (14,8.10) (6,6.13) (4,3.10) (12,9.13) (7,7.26) (5,4.74)};
\addplot[densityred, thick, domain=2:20] {3+0.5*x};
\end{axis}
\begin{axis}[an, at={(0,-5.1cm)}, title={III: a line plus one outlier}]
\addplot[only marks, mark=*, mark size=1.3pt, sourceblue] coordinates {(10,7.46) (8,6.77) (13,12.74) (9,7.11) (11,7.81) (14,8.84) (6,6.08) (4,5.39) (12,8.15) (7,6.42) (5,5.73)};
\addplot[densityred, thick, domain=2:20] {3+0.5*x};
\draw[trapred, thick] (axis cs:13,12.74) circle[radius=0.17cm];
\node[font=\scriptsize, trapred, anchor=west] at (axis cs:13.8,12.74) {outlier};
\end{axis}
\begin{axis}[an, at={(6.3cm,-5.1cm)}, title={IV: one point does all the work}]
\addplot[only marks, mark=*, mark size=1.3pt, sourceblue] coordinates {(7.65,5.25) (8.35,5.56) (7.65,5.76) (8.35,6.58) (7.65,6.89) (8.35,7.04) (7.65,7.71) (8.35,7.91) (7.65,8.47) (8.35,8.84)};
\addplot[only marks, mark=*, mark size=1.8pt, trapred] coordinates {(19,12.50)};
\addplot[densityred, thick, domain=2:20] {3+0.5*x};
\node[font=\scriptsize, anchor=north] at (axis cs:8.5,4.2) {10 points, all at $x=8$};
\node[font=\scriptsize, trapred, align=right, anchor=east] at (axis cs:18.4,12.9) {one point at $x=19$};
\end{axis}
\end{tikzpicture}
\caption{Anscombe's quartet. Same means, variances, correlation ($0.816$) and
fitted line (red). Only panel I is described well by ``linear with $r = 0.82$''.
II is a smooth curve, III is a perfect line ruined by one point, and IV has no
relationship at all apart from one point. (In IV the ten points at $x = 8$ are
nudged alternately left and right so that all ten can be seen.)}
\label{fig:two-variables-anscombe}
\end{figure}

Each panel teaches a separate lesson. Panel II: $r$ only measures fit to a line,
and a curve that fits perfectly gets $0.82$. Panel III: one outlier dragged $r$
down from $1$ to $0.82$. Panel IV: one \textbf{high-leverage} point, a point whose
$x$ value sits far from all the others, manufactured a correlation out of nothing; remove it and $x$ has no variation left at all.

\section{When $r$ misses the relationship}
\label{sec:corr-nonlinear}

\subsection{A perfect relationship with $r = 0$}

Take $x = -3, -2, -1, 0, 1, 2, 3$ and $y = x^2$, so $y = 9,4,1,0,1,4,9$. Here $y$
is completely determined by $x$. Yet $\xbar = 0$, $\ybar = 28/7 = 4$, and the
products $(x_i - \xbar)(y_i - \ybar) = x_i(y_i - 4)$ are
\[
  -15,\ 0,\ 3,\ 0,\ -3,\ 0,\ 15,
\]
which sum to $0$. So $r = 0$.

\begin{figure}[H]
\centering
\begin{tikzpicture}
\begin{axis}[width=9cm, height=6cm, xmin=-3.6, xmax=3.6, ymin=-1.5, ymax=10.5,
  xlabel={$x$}, ylabel={$y = x^2$}, xtick={-3,...,3}]
\fill[targetgreen, opacity=0.2] (axis cs:0,4) rectangle (axis cs:3,9);
\fill[targetgreen, opacity=0.2] (axis cs:0,4) rectangle (axis cs:-1,1);
\fill[trapred, opacity=0.2] (axis cs:0,4) rectangle (axis cs:-3,9);
\fill[trapred, opacity=0.2] (axis cs:0,4) rectangle (axis cs:1,1);
\addplot[gray, domain=-3.3:3.3, samples=60] {x^2};
\draw[dashed, gray] (axis cs:0,-1.5) -- (axis cs:0,10.5);
\draw[dashed, gray] (axis cs:-3.6,4) -- (axis cs:3.6,4);
\addplot[only marks, mark=*, mark size=2pt, sourceblue] coordinates {(-3,9) (-2,4) (-1,1) (0,0) (1,1) (2,4) (3,9)};
\node[font=\small] at (axis cs:-1.5,6.5) {area $-15$};
\node[font=\small] at (axis cs:1.5,6.5) {area $+15$};
\node[font=\scriptsize] at (axis cs:-0.5,2.5) {$+3$};
\node[font=\scriptsize] at (axis cs:0.5,2.5) {$-3$};
\node[font=\scriptsize, gray!50!black, anchor=south east] at (axis cs:3.35,4.05) {$\ybar = 4$};
\node[font=\small, align=center, anchor=south] at (axis cs:2.0,-1.4) {areas sum to $0$,\\ so $r = 0$};
\end{axis}
\end{tikzpicture}
\caption{Points exactly on a parabola, with each point's rectangle to the point of
means $(0, 4)$. The three points on the dashed lines, $(-2,4)$, $(0,0)$ and
$(2,4)$, have zero area. Each rectangle on the right cancels its mirror image on
the left, so the covariance and $r$ are exactly $0$. Zero correlation means no
\emph{linear} trend, not no relationship.}
\label{fig:two-variables-parabola}
\end{figure}

The lesson runs in one direction. If $x$ and $y$ are independent (knowing one tells
you nothing about the other), then in a large sample $r$ is near $0$. The converse fails: $r=0$ does
not imply independence.

\subsection{Spearman rank correlation}

Sometimes the relationship is curved but still always increasing. Suppose a service
is tested at load levels $x = 1, 2, \dots, 8$ and the queue length is
$y = 2, 4, 8, \dots, 256$: it doubles with every step. Bigger $x$ always means
bigger $y$, yet $r \approx 0.850$, because the points bend away from any line.

A relationship where $y$ always goes up when $x$ goes up (or always down) is
called \textbf{monotone}. To measure monotone association, replace each value by
its \textbf{rank}, its position when the values are sorted from smallest (rank 1)
to largest (rank $n$), and compute Pearson's $r$ on the ranks.

\begin{definition}[Spearman rank correlation]
The \textbf{Spearman correlation} $r_s$ is the Pearson correlation of the ranks of
the $x$ values with the ranks of the $y$ values. Tied values each get the average
of the ranks they occupy. When there are no ties,
\[
  r_s = 1 - \frac{6\sum_{i=1}^n d_i^2}{n(n^2-1)},
\]
where $d_i$ is the difference between the rank of $x_i$ and the rank of $y_i$.
\end{definition}

For the doubling queue, the ranks of $y$ are $1, 2, \dots, 8$, identical to the
ranks of $x$. Every $d_i = 0$ and $r_s = 1$.

\begin{figure}[H]
\centering
\begin{tikzpicture}
\begin{axis}[name=raw, scale only axis, width=4.3cm, height=3.4cm,
  xmin=0, xmax=9, ymin=-15, ymax=275, xlabel={load level $x$}, ylabel={queue length $y$},
  title={raw values: $r \approx 0.85$}, title style={font=\small},
  tick label style={font=\scriptsize}]
\addplot[only marks, mark=*, mark size=1.8pt, sourceblue] coordinates {(1,2) (2,4) (3,8) (4,16) (5,32) (6,64) (7,128) (8,256)};
\end{axis}
\begin{axis}[at={(6.6cm,0)}, scale only axis, width=4.3cm, height=3.4cm,
  xmin=0, xmax=9, ymin=0, ymax=9, xlabel={rank of $x$}, ylabel={rank of $y$},
  title={ranks: $r_s = 1$}, title style={font=\small},
  tick label style={font=\scriptsize}]
\addplot[only marks, mark=*, mark size=1.8pt, quotientorange] coordinates {(1,1) (2,2) (3,3) (4,4) (5,5) (6,6) (7,7) (8,8)};
\addplot[gray, dashed, domain=0.5:8.5] {x};
\end{axis}
\end{tikzpicture}
\caption{Ranking straightens any monotone relationship. The doubling curve on the
left becomes a perfect line on the right, so Spearman's $r_s$ is 1 while Pearson's
$r$ is only 0.85.}
\label{fig:two-variables-spearman}
\end{figure}

\begin{example}[Spearman by hand]
For the six requests the $x$ values are already in order, with ranks
$1,2,3,4,5,6$. The latencies $12, 19, 13, 20, 22, 16$ have ranks $1, 4, 2, 5, 6, 3$.
The differences are $d = 0, -2, 1, -1, -1, 3$, so $\sum d_i^2 = 0+4+1+1+1+9 = 16$
and
\[
  r_s = 1 - \frac{6 \times 16}{6 \times 35} = 1 - \frac{96}{210} \approx 0.543 .
\]
This is close to Pearson's $0.528$, as expected for data without strong curvature.
\end{example}

Ranks also tame outliers. A rank can be at most $n$, however extreme the value, so a
single wild point cannot dominate. In the outlier panel of
Figure~\ref{fig:two-variables-gallery}, Pearson's $r$ is $0.729$ with the outlier
and $0.960$ without; Spearman's $r_s$ is $0.801$ with it. The same logic made the
median more robust than the mean in Chapter 2.

\begin{lstlisting}[language=Python]
from scipy import stats
x = [1, 2, 3, 4, 5, 6, 7, 8]; y = [2**k for k in x]
stats.pearsonr(x, y).statistic    # 0.850
stats.spearmanr(x, y).statistic   # 1.0
\end{lstlisting}

\section{Correlation is not causation}
\label{sec:corr-causation}

Across the days of a year, a city's ice cream sales and its number of drownings are
positively correlated. Ice cream does not cause drowning. Hot weather drives both:
more people buy ice cream and more people swim. Temperature is a
\textbf{confounder}, a third variable that influences both the explanatory variable
and the response and so creates an association between them even when neither
affects the other.

A \textbf{causal diagram} draws each variable as a box and each direct causal
influence as an arrow. It is a picture of the story we believe generated the data.
An observed correlation between $X$ and $Y$ is consistent with several different
diagrams.

\begin{figure}[H]
\centering
\begin{tikzpicture}[
  v/.style={draw, rounded corners, fill=gray!8, minimum width=1.3cm, minimum height=0.7cm, font=\small},
  ar/.style={-{Stealth[length=2.4mm]}, thick}]
\node[v] (x1) at (0,0) {$X$};
\node[v] (y1) at (2.2,0) {$Y$};
\draw[ar, sourceblue] (x1) -- (y1);
\node[font=\small, align=center] at (1.1,-1.1) {(a) $X$ causes $Y$\\ \scriptsize e.g.\ payload $\to$ latency};

\node[v] (x2) at (4.7,0) {$X$};
\node[v] (y2) at (6.9,0) {$Y$};
\draw[ar, quotientorange] (y2) -- (x2);
\node[font=\small, align=center] at (5.8,-1.1) {(b) $Y$ causes $X$\\ \scriptsize e.g.\ crashes $\to$ more RAM};

\node[v] (x3) at (9.6,0) {$X$};
\node[v] (y3) at (11.8,0) {$Y$};
\node[v] (z3) at (10.7,1.5) {$Z$};
\draw[ar, trapred] (z3) -- (x3);
\draw[ar, trapred] (z3) -- (y3);
\draw[dashed, gray] (x3) -- (y3);
\node[font=\small, align=center] at (10.7,-1.1) {(c) $Z$ confounds\\ \scriptsize e.g.\ heat $\to$ ice cream, drownings};
\end{tikzpicture}
\caption{Three causal stories that all produce a correlation between $X$ and $Y$.
Arrows are causes; the examples underneath come from this chapter. Data on $X$
and $Y$ alone cannot tell the three apart. In (c) the dashed line is an
association with no arrow behind it: $X$ and $Y$ move together only because $Z$
pushes both.}
\label{fig:two-variables-dags}
\end{figure}

Concretely:
\begin{itemize}
  \item \textbf{Reverse causation.} Servers with more RAM have more crashes in a
    fleet. Perhaps RAM causes crashes, but more likely the team upgraded the RAM
    on servers that were already crashing.
  \item \textbf{Confounding.} Users who enable dark mode have higher retention. Dark
    mode is buried in settings, so the users who find it are engaged power users;
    engagement drives both.
  \item \textbf{Chance.} Search enough pairs of variables and some will correlate
    strongly by luck. Chapter 11 treats this as the multiple testing problem.
\end{itemize}

\begin{intuitionbox}
The way to break the confounder's grip is to set $X$ yourself, at random. If a coin
flip decides who gets dark mode, then engagement cannot influence who has it, so
arrow $Z \to X$ is cut. Any remaining difference in retention must flow along
$X \to Y$, up to chance. This is why randomized experiments (A/B tests) are the
gold standard for causal claims. Chapter 5 develops the design of studies and
experiments.
\end{intuitionbox}

\begin{historybox}
Francis Galton, a Victorian scientist and cousin of Charles Darwin, spent the 1870s
and 1880s measuring parents and children: heights, arm lengths, sweet-pea seed
sizes. He was stuck on how to put a single number on ``how strongly does a child's
height follow a parent's''. In 1888 he proposed measuring both variables in units
of their own spread, which is the z-score idea, and called the resulting slope the
``co-relation''. Karl Pearson turned this into the formula
$r = s_{xy}/(s_x s_y)$ in 1896, which is why it carries his name. Eight years later
the psychologist Charles Spearman introduced the rank version to study
intelligence test scores, where only the ordering seemed meaningful. Pearson's
colleague George Udny Yule noticed as early as 1903 that combining groups could
reverse an association, the effect now named after Edward Simpson, who wrote it up
in 1951. Galton's original goal, predicting one variable from another, became
regression (Chapter 12).
\end{historybox}

\section{Simpson's paradox}
\label{sec:corr-simpson}

Confounding can do more than create an association. It can flip one. A 1986 study
in the \emph{British Medical Journal} compared two treatments for kidney stones:
open surgery (A) and a less invasive procedure through a small puncture (B). A
treatment counted as a success if, three months later, the stone was gone or only
tiny fragments remained.

\begin{center}
\begin{tabular}{lccc}
\toprule
 & Small stones & Large stones & All patients \\
\midrule
Treatment A (open surgery) & $81/87 = 93\%$ & $192/263 = 73\%$ & $273/350 = 78\%$ \\
Treatment B (puncture) & $234/270 = 87\%$ & $55/80 = 69\%$ & $289/350 = 83\%$ \\
\bottomrule
\end{tabular}
\end{center}

Read the table one column at a time. For small stones, A wins: 93\% versus 87\%.
For large stones, A wins: 73\% versus 69\%. For all patients together, B wins:
83\% versus 78\%. Every patient has either a small or a large stone, so it seems
impossible that A is better in both groups and worse overall. It is not
impossible, and the arithmetic is right.

\begin{figure}[H]
\centering
\begin{tikzpicture}
\begin{axis}[ybar, width=11cm, height=6.5cm, bar width=0.7cm,
  symbolic x coords={Small stones, Large stones, All patients}, xtick=data,
  ymin=0, ymax=112, ytick={0,20,40,60,80,100}, ylabel={success rate (\%)},
  clip=false,
  enlarge x limits=0.25,
  nodes near coords, nodes near coords style={font=\scriptsize, /pgf/number format/fixed, /pgf/number format/precision=1, /pgf/number format/fixed zerofill},
  legend style={at={(0.5,1.03)}, anchor=south, legend columns=2, font=\small, draw=none}, legend cell align=left]
\addplot[fill=sourceblue, draw=sourceblue!70!black] coordinates {(Small stones,93.1) (Large stones,73.0) (All patients,78.0)};
\addplot[fill=quotientorange, draw=quotientorange!70!black] coordinates {(Small stones,86.7) (Large stones,68.8) (All patients,82.6)};
\legend{A (open surgery), B (puncture)}
\node[font=\small\bfseries, sourceblue!80!black] at (axis cs:Small stones,106) {A wins};
\node[font=\small\bfseries, sourceblue!80!black] at (axis cs:Large stones,86) {A wins};
\node[font=\small\bfseries, quotientorange!80!black] at (axis cs:All patients,96) {B wins};
\end{axis}
\end{tikzpicture}
\caption{Simpson's paradox in the kidney-stone data (bars start at 0). Compare
the two bars in each group: A is taller within small stones and within large
stones, yet B is taller when the groups are combined.}
\label{fig:two-variables-simpson-bars}
\end{figure}

The resolution is in who got which treatment. Doctors gave the invasive surgery to
the hard cases: 263 of A's 350 patients (75\%) had large stones, but only 80 of
B's 350 (23\%) did. Large stones are harder to treat by any method. So A's overall
rate is mostly a large-stone rate, and B's is mostly a small-stone rate. Each
overall rate is a weighted average of the two group rates:
\begin{align*}
  \text{A overall} &= \tfrac{87}{350}(93.1\%) + \tfrac{263}{350}(73.0\%) \approx 78.0\%,\\
  \text{B overall} &= \tfrac{270}{350}(86.7\%) + \tfrac{80}{350}(68.8\%) \approx 82.6\%.
\end{align*}
The weights differ, and the weights carry the stone-size effect.

\begin{figure}[H]
\centering
\begin{tikzpicture}
\begin{axis}[width=10.5cm, height=7cm, xmin=0, xmax=1, ymin=0.6, ymax=1.0,
  xlabel={fraction of patients with large stones},
  ylabel={success rate},
  xtick={0,0.229,0.5,0.751,1}, xticklabels={0,0.23,0.5,0.75,1},
  ytick={0.6,0.7,0.8,0.9,1.0},
  legend style={at={(0.02,0.03)}, anchor=south west, font=\small}, legend cell align=left]
\addplot[sourceblue, very thick, domain=0:1] {0.931 + (0.730-0.931)*x};
\addplot[quotientorange, very thick, domain=0:1] {0.867 + (0.6875-0.867)*x};
\draw[densely dotted, gray] (axis cs:0.751,0.6) -- (axis cs:0.751,0.780);
\draw[densely dotted, gray] (axis cs:0.229,0.6) -- (axis cs:0.229,0.826);
\addplot[only marks, mark=*, mark size=3.5pt, sourceblue] coordinates {(0.751,0.780)};
\addplot[only marks, mark=*, mark size=3.5pt, quotientorange] coordinates {(0.229,0.826)};
\draw[gray] (axis cs:0.751,0.780) -- (axis cs:0.80,0.895);
\node[font=\small, anchor=south] at (axis cs:0.80,0.895) {A's actual mix: 78\%};
\draw[gray] (axis cs:0.229,0.826) -- (axis cs:0.42,0.955);
\node[font=\small, anchor=south] at (axis cs:0.42,0.955) {B's actual mix: 83\%};
\legend{A (open surgery), B (puncture)}
\end{axis}
\end{tikzpicture}
\caption{Each line shows a treatment's overall success rate if its patients had a
given fraction of large stones: the left end is its small-stone rate, the right
end its large-stone rate. A's line is above B's everywhere, so A is better at
every fixed mix. The dots mark the actual mixes: A treated mostly large stones
(75\%), B mostly small (23\% large), and that pushes A's dot below B's.}
\label{fig:two-variables-simpson-lines}
\end{figure}

\begin{figure}[H]
\centering
\begin{tikzpicture}[
  v/.style={draw, rounded corners, fill=gray!8, minimum height=0.75cm, font=\small, align=center},
  ar/.style={-{Stealth[length=2.4mm]}, thick}]
\node[v] (z) at (0,1.6) {stone size};
\node[v] (x) at (-2.6,0) {treatment};
\node[v] (y) at (2.6,0) {success};
\draw[ar, trapred] (z) -- node[above left, font=\scriptsize] {doctors choose} (x);
\draw[ar, trapred] (z) -- node[above right, font=\scriptsize] {large is harder} (y);
\draw[ar, sourceblue] (x) -- node[below, font=\scriptsize] {effect we want} (y);
\end{tikzpicture}
\caption{Causal diagram for the kidney-stone study. Stone size affects both which
treatment a patient gets and whether it succeeds, so it confounds the comparison.
Comparing within stone size removes its influence.}
\label{fig:two-variables-simpson-dag}
\end{figure}

\begin{trapbox}
There is no rule ``always trust the split table''. Which table answers your question
depends on the causal story. Here stone size is measured \emph{before} treatment and
affects it, so you compare within stone size, and A is the better treatment. If the
splitting variable were something the treatment itself caused (say, blood pressure
measured after surgery), splitting would block part of the treatment's effect and
the combined table would be the right one. The numbers alone cannot tell you;
the diagram can.
\end{trapbox}

The most famous case involves university admissions. In fall 1973 the University
of California, Berkeley admitted about 44\% of its 8442 male graduate applicants
and about 35\% of its 4321 female applicants. A 1975 analysis by Bickel, Hammel and
O'Connell looked department by department and found that in most departments women
were admitted at a similar or higher rate than men. Women had applied
disproportionately to departments that admitted few applicants of either sex.
The department chosen played the role of stone size.

\begin{recoverybox}
Check your understanding: could A have won in both groups and \emph{also} overall?
(Yes, if both treatments had been given to similar mixes of patients. The reversal
needs the groups to differ in both size mix and baseline difficulty.)
\end{recoverybox}

\section{Ecological fallacy}
\label{sec:corr-ecological}

A correlation computed from group averages (countries, states, schools, server
racks) is called an \textbf{ecological correlation}. The \textbf{ecological
fallacy} is reading it as if it described individuals.

Averaging removes the scatter within each group. What is left is usually a much
tighter pattern than the individuals show, and sometimes a pattern in a different
direction.

\begin{figure}[H]
\centering
\begin{tikzpicture}
\begin{axis}[width=10.5cm, height=7cm, xmin=-2.5, xmax=10.5, ymin=-3, ymax=13,
  xlabel={$x$}, ylabel={$y$},
  legend style={at={(0.5,1.03)}, anchor=south, legend columns=6, font=\scriptsize, draw=none}, legend cell align=left]
\addplot[only marks, mark=*, mark size=1.3pt, sourceblue!70] coordinates {(5.43,2.64) (-1.92,1.79) (2.84,0.94) (1.26,2.40) (1.44,4.32) (1.82,2.74) (-1.06,5.37) (1.80,2.82) (0.78,3.31) (7.48,6.66) (2.53,4.46) (1.60,2.15)};
\addplot[only marks, mark=*, mark size=1.3pt, targetgreen!80] coordinates {(2.70,-2.10) (3.85,6.37) (6.09,2.01) (2.56,0.45) (2.60,4.73) (4.59,5.66) (1.57,3.14) (2.52,1.75) (4.40,4.18) (3.92,4.01) (3.14,7.22) (4.06,5.77)};
\addplot[only marks, mark=*, mark size=1.3pt, quotientorange!80] coordinates {(5.41,2.96) (6.88,6.70) (4.77,4.44) (3.62,8.23) (6.03,2.09) (6.03,4.55) (4.75,8.13) (3.85,9.11) (5.47,1.82) (1.11,5.47) (6.20,7.29) (5.88,5.22)};
\addplot[only marks, mark=*, mark size=1.3pt, structurepurple!70] coordinates {(8.60,7.40) (4.15,5.22) (6.62,4.87) (4.38,4.87) (8.31,9.65) (6.02,3.42) (5.46,5.32) (3.31,5.92) (8.75,7.13) (7.26,7.90) (7.26,3.88) (7.88,2.82)};
\addplot[only marks, mark=*, mark size=1.3pt, trapred!70] coordinates {(7.66,7.23) (9.61,12.35) (8.88,10.04) (8.01,10.03) (7.16,2.25) (8.14,5.99) (7.28,5.40) (8.97,7.91) (6.39,1.72) (7.88,9.32) (7.49,9.09) (8.54,3.87)};
\addplot[only marks, mark=diamond*, mark size=4.5pt, black] coordinates {(2.00,3.30) (3.50,3.60) (5.00,5.50) (6.50,5.70) (8.00,7.10)};
\legend{group 1, group 2, group 3, group 4, group 5, group means}
\node[font=\small, align=left, anchor=north west, fill=white, inner sep=2pt] at (axis cs:-2.3,12.7) {60 individuals (dots): $r \approx 0.50$\\ 5 group means (diamonds): $r \approx 0.97$};
\end{axis}
\end{tikzpicture}
\caption{Sixty individuals in five groups (colours) and the five group means
(black diamonds). The diamonds lie almost on a line ($r \approx 0.97$), but the
individuals around them scatter widely ($r \approx 0.50$). A correlation of
averages overstates how well $x$ predicts $y$ for any one person.}
\label{fig:two-variables-ecological}
\end{figure}

The classic real example is W.\,S.\ Robinson's 1950 study of the 1930 US census.
Across the 48 states, the percentage of residents born abroad correlated
\emph{negatively} with the percentage who were illiterate: $r \approx -0.53$.
Read naively, immigrants were more literate. Across individuals the correlation was
\emph{positive}, about $0.12$: immigrants were somewhat more likely to be
illiterate. The state-level pattern arose because immigrants tended to settle in
states whose native-born population had high literacy.
Figure~\ref{fig:two-variables-ecological-reversal} shows how such a reversal can
happen.

\begin{figure}[H]
\centering
\begin{tikzpicture}
\begin{axis}[width=10.5cm, height=6.8cm, xmin=0, xmax=10, ymin=0.5, ymax=11,
  xlabel={$x$}, ylabel={$y$},
  legend style={at={(0.98,0.98)}, anchor=north east, font=\scriptsize}, legend cell align=left]
\addplot[only marks, mark=*, mark size=1.6pt, sourceblue] coordinates {(0.75,8.30) (1.25,8.15) (1.75,8.90) (2.25,8.90) (2.75,9.85) (3.25,9.90)};
\addplot[only marks, mark=*, mark size=1.6pt, targetgreen!80!black] coordinates {(2.75,5.90) (3.25,6.70) (3.75,6.55) (4.25,7.30) (4.75,7.30) (5.25,8.25)};
\addplot[only marks, mark=*, mark size=1.6pt, quotientorange] coordinates {(4.75,4.25) (5.25,4.30) (5.75,5.10) (6.25,4.95) (6.75,5.70) (7.25,5.70)};
\addplot[only marks, mark=*, mark size=1.6pt, structurepurple] coordinates {(6.75,1.70) (7.25,2.65) (7.75,2.70) (8.25,3.50) (8.75,3.35) (9.25,4.10)};
\addplot[sourceblue, thick, domain=0.6:3.4] {9 + 0.8*(x-2)};
\addplot[targetgreen!80!black, thick, domain=2.6:5.4] {7 + 0.8*(x-4)};
\addplot[quotientorange, thick, domain=4.6:7.4] {5 + 0.8*(x-6)};
\addplot[structurepurple, thick, domain=6.6:9.4] {3 + 0.8*(x-8)};
\addplot[black, densely dashed, thick, domain=1.4:8.6] {11 - x};
\addplot[only marks, mark=diamond*, mark size=4.5pt, black] coordinates {(2,9) (4,7) (6,5) (8,3)};
\legend{group 1, group 2, group 3, group 4}
\node[font=\small, align=left, anchor=south west] at (axis cs:0.15,0.7) {inside each group\\ (solid lines): $r \approx +0.95$\\[2pt] across the 4 group\\ means (dashed): $r = -1$};
\end{axis}
\end{tikzpicture}
\caption{An ecological reversal, built to mirror Robinson's census data. Inside
every group, $y$ rises with $x$ (solid lines). But the groups with larger $x$ sit
lower, so the group means (diamonds) fall along the dashed line. The group-level
correlation has the opposite sign from the one that holds for every individual
within a group.}
\label{fig:two-variables-ecological-reversal}
\end{figure}

\begin{trapbox}
``Countries that eat more chocolate win more Nobel prizes'' is a statement about
countries. It says nothing about whether a chocolate-eating person is more likely
to win one. Before using a correlation, ask what one data point is: a person, a
request, a server, a city, a week.
\end{trapbox}

\section{Restriction of range}
\label{sec:corr-range}

A university admits students using an entrance exam score and later finds that,
among admitted students, the exam barely correlates with first-year grades. Is the
exam useless? Not necessarily.

\begin{figure}[H]
\centering
\begin{tikzpicture}
\begin{axis}[width=11cm, height=7cm, xmin=190, xmax=780, ymin=1.6, ymax=4.3,
  xlabel={entrance score}, ylabel={first-year GPA},
  legend style={at={(0.02,0.98)}, anchor=north west, font=\small}, legend cell align=left]
\fill[targetgreen, opacity=0.12] (axis cs:600,1.6) rectangle (axis cs:780,4.3);
\addplot[only marks, mark=*, mark size=1.2pt, gray!60] coordinates {(552,3.64) (316,3.11) (507,2.79) (424,2.49) (488,3.26) (577,3.42) (326,2.09) (400,3.19) (581,3.30) (422,2.84) (496,3.30) (526,3.60) (427,3.00) (567,2.92) (382,3.06) (569,3.15) (556,2.95) (539,2.93) (533,3.43) (475,3.24) (399,2.68) (362,2.65) (521,3.02) (483,3.41) (432,2.80) (525,2.76) (433,2.89) (436,2.66) (586,3.54) (506,3.05) (510,2.91) (456,2.59) (584,3.21) (497,3.09) (571,3.66) (482,2.86) (534,2.85) (589,3.21) (594,3.60) (567,3.11) (327,2.90) (473,2.69) (356,2.64) (483,3.30) (304,2.41) (497,3.47) (393,1.82) (360,2.76) (556,2.91) (490,3.01) (400,2.76) (505,2.78) (416,2.94) (588,3.65) (293,2.03) (435,2.48) (404,2.81) (371,2.56) (288,2.19) (548,3.38) (458,3.11) (536,3.34) (442,2.43) (400,2.60) (571,3.14) (468,2.89) (447,3.17) (569,2.93) (500,2.80) (434,2.99) (472,2.94) (519,2.90) (559,3.32) (559,3.53) (482,2.51) (459,2.99) (515,3.25) (500,2.71) (393,2.41) (467,3.44) (597,2.73) (499,3.11) (581,3.34) (442,3.01) (451,2.19) (434,2.65) (397,2.83) (522,2.63) (286,2.46) (527,2.91) (214,2.30) (449,2.62) (477,2.61) (574,3.22) (451,3.00) (509,2.57) (496,2.87) (321,2.37) (559,2.98) (479,3.15) (413,2.68) (361,2.25) (535,2.98) (539,3.28) (428,2.79) (483,2.77) (420,2.93) (497,3.13) (390,2.48) (484,3.03) (464,3.02) (402,2.72) (423,3.11) (561,3.07) (445,2.99) (477,3.01) (514,3.17) (402,2.59) (281,2.50) (443,3.17) (497,2.42) (435,2.87) (510,3.25) (568,3.15)};
\addplot[only marks, mark=*, mark size=1.6pt, targetgreen!70!black] coordinates {(683,3.80) (752,4.05) (669,3.06) (641,2.83) (652,3.71) (610,3.41) (601,3.84) (610,3.30) (653,3.12) (659,3.66) (660,3.26) (633,3.18) (679,3.09) (606,3.47) (601,3.24) (626,3.26) (617,3.66) (667,3.99) (732,3.50) (622,3.66) (616,3.01) (648,3.84) (620,3.09) (630,2.70) (641,2.93) (663,3.25)};
\draw[dashed, targetgreen!60!black, thick] (axis cs:600,1.6) -- (axis cs:600,4.3);
\node[font=\small, anchor=south east, fill=white, inner sep=2pt] at (axis cs:775,1.65) {admitted only: $r \approx 0.30$};
\node[font=\small, anchor=north west] at (axis cs:197,3.72) {all 150 together: $r \approx 0.70$};
\legend{not admitted (score $< 600$), admitted (score $\ge 600$)}
\end{axis}
\end{tikzpicture}
\caption{Restriction of range. Over all 150 applicants (grey and green dots
together) the score and GPA correlate at $0.70$. Keep only those with a score of at least 600 and the correlation falls to
about $0.30$. The relationship did not change; the range of $x$ shrank.}
\label{fig:two-variables-range}
\end{figure}

Correlation compares the trend in $y$ to the scatter around it. If you keep only a
narrow slice of $x$ values, the trend across the slice is small, but the scatter
of $y$ around the trend is as large as ever. So $r$ drops. The same thing happens
when you study only the employees who passed the hiring screen, only the requests
that did not time out, or only the patients healthy enough to enter a trial.
Correlations from truncated samples cannot be compared directly to correlations
from full ones.

\begin{recoverybox}
The reverse also holds. If you sample only the very small and very large payloads
and skip the middle, would $r$ go up or down? (Up. A wider spread in $x$ makes the
same trend stand out more against the same scatter.)
\end{recoverybox}

\section{Preview: the regression line}
\label{sec:corr-regression-preview}

Correlation is symmetric and unitless. Often we want something asymmetric with
units: a rule that predicts $y$ from $x$. The natural candidate is a line through
the point of means $(\xbar, \ybar)$. Its slope turns out to be
\[
  b = r \,\frac{s_y}{s_x}.
\]
In words: when $x$ rises by one standard deviation, the predicted $y$ rises by $r$
standard deviations. For the six requests,
\[
  b = 0.528 \times \frac{4}{2.366} \approx 0.893 \text{ ms per KB},
\]
which is also $s_{xy}/s_x^2 = 5/5.6 = 25/28$. The intercept is
$\ybar - b\xbar = 17 - 0.893 \times 4 \approx 13.43$, so the line is
$\hat y \approx 13.43 + 0.893x$.

\begin{figure}[H]
\centering
\begin{tikzpicture}
\begin{axis}[width=10.5cm, height=7cm, xmin=0, xmax=8, ymin=8, ymax=26,
  xlabel={payload $x$ (KB)}, ylabel={latency $y$ (ms)},
  xtick={1,...,7},
  legend style={at={(0.02,0.98)}, anchor=north west, font=\small}, legend cell align=left]
\addplot[densityred, thick, domain=0:8] {13.4286 + 0.892857*x};
\addplot[gray, dashed, domain=0.8:7.2] {17 + 1.690309*(x-4)};
\addplot[only marks, mark=*, mark size=2pt, sourceblue] coordinates {(1,12) (2,19) (3,13) (5,20) (6,22) (7,16)};
\addplot[only marks, mark=x, mark size=4pt, black, thick] coordinates {(4,17)};
\draw[thick, black!70] (axis cs:4,17) -- (axis cs:6.366,17) -- (axis cs:6.366,19.113);
\node[font=\scriptsize, anchor=north] at (axis cs:5.18,17) {$s_x = 2.37$ KB};
\node[font=\scriptsize, align=left, anchor=west] at (axis cs:6.42,18.05) {$r s_y$\\ $= 2.11$ ms};
\node[font=\scriptsize, anchor=south east] at (axis cs:3.9,17.2) {$(\xbar,\ybar) = (4, 17)$};
\legend{regression line slope $r s_y/s_x$, SD line slope $s_y/s_x$}
\end{axis}
\end{tikzpicture}
\caption{The regression line (red) passes through $(\xbar,\ybar)$ (black cross).
The step drawn from the cross shows its slope: one standard deviation to the
right in $x$ ($2.37$ KB), it rises only $r$ standard deviations in $y$
($0.528 \times 4 = 2.11$ ms), so it is flatter than the dashed line that rises a full $s_y$ per $s_x$.
Chapter 12 shows that this line minimizes the sum of squared vertical errors.}
\label{fig:two-variables-regression}
\end{figure}

Because $\abs{r} \le 1$, the prediction for $y$ is never more standard deviations
from its mean than $x$ is from its mean, and strictly fewer unless the points lie
on a line. Tall parents tend to have children who are
tall, but less extreme. This is \emph{regression to the mean}, the effect Galton was
studying, and it is pure arithmetic of the slope formula, not a force of nature.

\begin{connectionbox}
In machine learning, the correlation matrix of features is a standard first look at
a data set, and features with $\abs{r}$ near 1 are flagged as redundant (Chapter 13
calls this collinearity). Two cautions from this chapter carry over directly:
near-zero $r$ does not mean a feature is useless to a nonlinear model such as a tree
or a neural network, and a high $r$ between a feature and the label does not mean
the feature causes the label, so it may stop working when the world shifts.
In probability language, $r$ estimates
$\Corr(X,Y) = \Cov(X,Y)/(\SD(X)\SD(Y))$, where
$\Cov(X,Y) = \E[(X-\E X)(Y-\E Y)]$ (Probability From the Ground Up treats these
for random variables).
\end{connectionbox}

\section{Summary}
\label{sec:corr-summary}

\begin{itemize}
  \item \textbf{Two categorical variables}: make a contingency table, divide by row
    totals to get conditional distributions, and compare them with a segmented bar
    chart. No association means every row matches the marginal.
  \item \textbf{Two numerical variables}: draw a scatter plot; describe direction,
    form, strength and outliers.
  \item \textbf{Covariance} $s_{xy} = \frac{1}{n-1}\sum (x_i-\xbar)(y_i-\ybar)$ is the
    average signed rectangle area around $(\xbar,\ybar)$. Its size depends on units.
  \item \textbf{Pearson's} $r = \frac{1}{n-1}\sum z_{x,i}z_{y,i} = s_{xy}/(s_x s_y)$:
    unitless, symmetric, between $-1$ and $1$, equal to $\pm1$ only for a perfect line.
  \item $r$ measures \textbf{linear} association only. Anscombe's quartet: same $r$,
    four different pictures. A parabola can have $r=0$. Always plot.
  \item \textbf{Spearman's} $r_s$ is Pearson's $r$ on ranks; it measures monotone
    association and resists outliers.
  \item \textbf{Correlation is not causation}: reverse causation, confounders, and
    chance all create correlations. Randomization breaks confounding.
  \item \textbf{Simpson's paradox}: a confounder can reverse an association when groups
    are combined. The causal diagram decides which table to trust.
  \item \textbf{Ecological fallacy}: correlations of group averages are usually
    stronger than, and can even oppose, individual-level correlations.
  \item \textbf{Restriction of range} shrinks $r$ toward 0.
  \item \textbf{Regression slope} $b = r\,s_y/s_x$, line through $(\xbar,\ybar)$
    (Chapter 12).
\end{itemize}

Every trap in this chapter came from how the data were gathered: who chose the
treatment, which groups were averaged, which applicants were admitted. Chapter 5
takes that question head on: how sampling and study design decide what a data set
can tell you.

\section{Exercises}
\label{sec:corr-exercises}

\begin{enumerate}
  \item A/B test results by device. Desktop: 400 users, 48 converted. Mobile: 600
    users, 42 converted. (a) Give the conversion rate for each device and overall.
    (b) Under no association between device and conversion, how many desktop
    conversions would you expect?
  \item For the points $(1,2), (2,1), (3,4), (4,3), (5,5)$ compute $\xbar$, $\ybar$,
    $s_{xy}$, $s_x$, $s_y$ and $r$.
  \item Temperatures $x$ in degrees Celsius have correlation $r = 0.6$ with ice cream
    sales $y$. What is the correlation if temperature is converted to Fahrenheit,
    $x' = 32 + 1.8x$? What if $x$ is replaced by $-x$?
  \item Compute Pearson's $r$ and Spearman's $r_s$ for $(1,1), (2,2), (3,3), (4,4),
    (5,20)$. Which better describes ``$y$ goes up whenever $x$ goes up''?
  \item Give data with five points on the curve $y = \abs{x}$ for which $r = 0$
    exactly. Explain using the quadrant picture.
  \item A study finds that cities with more police officers per capita have more
    crime. Draw two different causal diagrams consistent with this.
  \item Hospital P has a 90\% survival rate and hospital Q has 85\%. Describe how Q
    could still be the better hospital for every type of patient. What additional
    variable would you ask for?
  \item An engineering team computes $r = 0.1$ between interview score and job
    performance, using only people they hired. Why might the interview still be a
    good predictor among all applicants?
  \item Prove that if every $y_i = a + bx_i$ with $b < 0$, then $r = -1$, by writing
    $z_{y,i}$ in terms of $z_{x,i}$.
  \item For the six requests in Section~\ref{sec:corr-covariance}, use the regression
    line to predict the latency for a 4 KB and a 10 KB payload. Which prediction do
    you trust less, and why?
\end{enumerate}

\subsection*{Answers}

\begin{enumerate}
  \item (a) Desktop $48/400 = 12\%$, mobile $42/600 = 7\%$, overall $90/1000 = 9\%$.
    (b) $400 \times 90/1000 = 36$.
  \item $\xbar = \ybar = 3$. Deviations $x$: $-2,-1,0,1,2$; $y$: $-1,-2,1,0,2$.
    Products $2,2,0,0,4$, sum $8$, so $s_{xy} = 8/4 = 2$. $s_x = s_y = \sqrt{10/4}
    \approx 1.581$. $r = 2/2.5 = 0.8$.
  \item Still $0.6$ (positive rescaling). With $-x$ it becomes $-0.6$.
  \item With $\xbar = 3$ and $\ybar = 6$, the products of deviations sum to $40$,
    $\sum(x_i-\xbar)^2 = 10$ and $\sum(y_i-\ybar)^2 = 250$, so
    $r = 40/\sqrt{2500} = 0.8$. The ranks of $y$ are $1,2,3,4,5$, so $r_s = 1$. Spearman
    describes the monotone pattern; Pearson is pulled around by the point $(5,20)$.
  \item $x = -2,-1,0,1,2$, $y = 2,1,0,1,2$. $\ybar = 1.2$; the products
    $x_i(y_i - 1.2)$ are $-1.6, 0.2, 0, -0.2, 1.6$, summing to $0$. Mirror-image
    points sit in opposite-sign quadrants and cancel.
  \item (i) Crime $\to$ more police (cities hire in response: reverse causation).
    (ii) City size or density $\to$ both police and crime (confounding).
  \item If Q treats a much larger share of high-risk patients, each risk group can
    have higher survival at Q while Q's overall rate is lower (Simpson's paradox).
    Ask for patient severity and compare within severity groups.
  \item Restriction of range: hires all scored above a cutoff, so the range of scores
    is narrow and $r$ is shrunk toward 0.
  \item $\ybar = a + b\xbar$ and $s_y = \abs{b}s_x = -b s_x$, so
    $z_{y,i} = b(x_i-\xbar)/(-b s_x) = -z_{x,i}$. Then
    $r = \frac{1}{n-1}\sum (-z_{x,i}^2) = -\frac{n-1}{n-1} = -1$.
  \item $\hat y(4) \approx 13.43 + 3.57 = 17.0$ ms (the mean, since $\xbar = 4$);
    $\hat y(10) \approx 13.43 + 8.93 = 22.4$ ms. The 10 KB prediction is less
    trustworthy: it lies outside the observed range 1--7 KB (extrapolation), and
    $r \approx 0.53$ from six points is a weak basis.
\end{enumerate}
